%% 3. Présentation de la mission  (FACULTATIF)
%% À garder si l'introduction ne suffit pas à poser le sujet ;
%% sinon fusionner avec l'introduction et retirer ce \input dans main.tex.
\chapter{Présentation de la mission}\label{ch:mission}

% RÔLE DE CE CHAPITRE : présenter les missions — le QUOI et le POURQUOI
% (problématique, besoin métier, objectifs, périmètre, contraintes).
%   - Les OUTILS et technologies mobilisés -> chap. Environnement technique.
%   - Les CHOIX entre technologies et l'APPRENTISSAGE -> chap. Veille.
%   - La RÉALISATION et les RÉSULTATS -> chap. Travaux effectués.
% On renvoie à ces chapitres par des \ref pour éviter les redites.

Au cours de mon alternance au sein de la division Identité \& Onboarding de
Namirial, j'ai contribué à deux missions principales, complétées par plusieurs
missions annexes. Les deux missions principales s'inscrivent dans le traitement
automatisé des pièces d'identité~: la première relève de la \textbf{lutte contre la
fraude documentaire}, la seconde de l'\textbf{extraction d'informations}. Les
missions annexes, plus ponctuelles, m'ont permis d'aborder d'autres facettes du
traitement intelligent de documents.

Cependant tous partagaient le même contexte des produits de vérification d'identité. Les modèles d'apprentissage automatique ne sont
pas une fin en soi, mais doivent être directement industrialisés dans la plateforme
et satisfaire des exigences fortes de fiabilité, de robustesse et de conformité.

\section{Première mission : classification colorimétrique selon les normes ICAO}

\subsection{Problématique et besoin métier}

Les pièces d'identité sont des documents fortement normés c'est-à-dire pour un type de document
et un pays donnés, la nature de la photo d'identité du titulaire est imposée par les
spécifications officielles. Celle-ci exigent selon les cas, une photo en couleur ou
en noir et blanc. Une photo ne respectant pas la règle attendue constitue
donc un indice de non-conformité, voire de falsification.

L'objectif de cette mission était de \textbf{déterminer automatiquement si la
photographie d'identité d'un document est en couleur ou en noir et blanc}, puis de
confronter ce résultat à la règle attendue pour le type de document concerné, afin
d'alimenter le dispositif de détection de fraude.

Le recours à un modèle d'apprentissage automatique se justifie par les limites d'une
approche purement déterministe. Une simple analyse des canaux chromatiques se révèle
insuffisante dans les cas ambigus. Par exemple, une photographie réellement en couleur, mais peu
contrastée et présentant un fond clair, pouvait être prise à tort pour une image en
noir et blanc. À l'inverse, un point de \textbf{vigilance éthique} a été identifié avec
le risque de \textbf{biais} conduisant à classer plus fréquemment en noir et blanc
les photographies de personnes à la peau foncée. Ces situations fréquentes
en conditions réelles (qualité de numérisation variable, éclairage, compression),
appelaient une décision plus fine qu'un seuillage ne peut offrir, et imposaient une
attention particulière à l'équité du modèle.

\subsection{Objectifs et périmètre}

Ce premier projet a été co-conduoit avec une collègue elle-aussi en alternance. Nous avons du le mené end-to-end. Il nous a permis de nous familiariser avec les outils utilisés par l'entreprise. Le périmètre incluait la constitution du jeu de
données, l'étude des caractéristiques discriminantes et l'entraînement du modèle~;
la \textbf{mise en industrialisation} du modèle, intervenue après ma contribution, en
était exclue, même si j'ai pu en suivre les enjeux.

La démarche de conception et les résultats sont détaillés au
chapitre~\ref{ch:travaux}, le choix de l'architecture au chapitre~\ref{ch:veille}, et
les outils internes employés au chapitre~\ref{ch:env}.

\subsection{Contraintes}

% À COMPLÉTER / À AFFINER. Pistes :
%   - Technique : exigence de légèreté et de rapidité du modèle pour une intégration
%     en production (motive le choix d'architecture, voir chap. Veille).
%   - Données : confidentialité des pièces d'identité (RGPD, anonymisation),
%     volume et déséquilibre éventuel des classes, hétérogénéité des acquisitions.
%   - Équité : maîtrise du biais évoqué ci-dessus.
Dans le cahier des charges on devait avoir un modèle très petit, très efficient et qui patate. Il ne devait pas y avoir de faux positif car on ne voulait pas que ce check de couleur soit bloquant.  

\section{Seconde mission : extraction généralisée d'informations}

Cette seconde mission portait sur une \textbf{approche généralisée de l'extraction
d'informations clés} (\emph{Key Information Extraction}) appliquée aux pièces
d'identité.

% À COMPLÉTER — partie majeure du mémoire (signalée comme telle).
% Au NIVEAU PRÉSENTATION (ici) :
%   - Besoin métier : quels champs extraire, depuis quels documents, et quel
%     processus existant (saisie manuelle ? OCR dédié par modèle de document ?)
%     cette approche « généralisée » vient remplacer.
%   - En quoi l'approche est « généralisée » (indépendante du type/pays de document)
%     et pourquoi c'est un enjeu (passage à l'échelle, maintenance).
%   - Objectifs, périmètre, contraintes (volumétrie, fiabilité, conformité).
% La conception, le code et les résultats iront au chap. Travaux (\ref{ch:travaux}).

\section{Missions annexes}

Trois interventions plus ponctuelles ont complété ces missions principales~; elles
sont présentées ici dans leurs grandes lignes et développées dans les chapitres
correspondants.

\paragraph{Évaluation de modèles d'OCR.}
À partir d'octobre~2025, une série de modèles d'OCR de nouvelle génération a été
publiée, annonçant de meilleures performances de lecture et une restitution
structurée du contenu (\emph{markdown}, tableaux). J'ai été chargée d'en évaluer
l'intérêt pour nos besoins. Cette démarche d'évaluation comparative, ainsi que ses
conclusions, relèvent de la veille technologique et sont traitées au
chapitre~\ref{ch:veille}.

\paragraph{Optimisation du Docprocessor.}
Le \emph{Docprocessor} est un service interne d'extraction d'information sur
documents (voir chap.~\ref{ch:env}). J'ai contribué à une tâche d'\textbf{optimisation
du temps de traitement} de son étape de classification. La solution mise en \oe uvre
et le gain obtenu sont décrits au chapitre~\ref{ch:travaux}.

\paragraph{Projet BUCAP.}
BUCAP est un projet client de \emph{débundling} de documents (séparation d'un
document composite en ses pièces constitutives). Une première version existait mais
demeurait fragile. La principale difficulté tenait à la \textbf{rareté des données}~:
seulement 21~exemples annotés, avec une vérité terrain (découpage des pages) très
incertaine, au point que les clients eux-mêmes divergeaient sur le découpage attendu.
Les pistes explorées pour pallier ce manque de données sont discutées au
chapitre~\ref{ch:veille}, et la réalisation au chapitre~\ref{ch:travaux}.

\section{Notions techniques nécessaires}

La compréhension des chapitres suivants suppose quelques notions de base~:

\begin{itemize}
  \item les \textbf{réseaux de neurones convolutifs} (CNN) pour la classification
        d'images, et en particulier les architectures \textbf{ResNet} (connexions
        résiduelles) et \textbf{ShuffleNet}, sa variante allégée~;
  \item les notions d'\textbf{ingénierie des caractéristiques} (\emph{feature
        engineering}) et de représentation de la couleur (espaces colorimétriques)~;
  \item les principes de l'\textbf{OCR} et de l'\textbf{extraction d'informations
        clés} (\emph{Key Information Extraction})~;
  \item l'usage de \textbf{grands modèles de langage} pour l'analyse de documents et
        la \textbf{contrainte de leur sortie} par un schéma typé ou une énumération~;
  \item les techniques d'\textbf{augmentation de données} (\emph{data augmentation}),
        dont l'\emph{inpainting}, face à un volume de données limité.
\end{itemize}
